\documentclass[aspectratio=169]{beamer}
\usepackage{beamerthemeAMU}
\usepackage{amsmath}
\usepackage{booktabs}
\usepackage{array}
\usepackage{listings}
\usepackage{hyperref}
\usepackage{tikz}
\usetikzlibrary{arrows.meta,positioning,calc,fit}

% TikZ styles for RAG diagrams (kept visually consistent with Week 1)
\colorlet{ragQueryBg}{AMU_acc5!35!AMU_lt1}
\colorlet{ragContextBg}{AMU_lt2!70!AMU_lt1}
\colorlet{ragLLMBg}{AMU_dk2!15!AMU_lt1}
\colorlet{ragStoreBg}{AMU_acc6!18!AMU_lt1}
\colorlet{ragArrow}{AMU_acc1}
\colorlet{ragGray}{AMU_dk1!65!AMU_lt1}

\tikzset{rag/box/.style={draw=AMU_dk1, rounded corners=6pt, inner sep=7pt, align=center, font=\small}}
\tikzset{rag/query/.style={rag/box, fill=ragQueryBg, minimum width=3.8cm}}
\tikzset{rag/llm/.style={rag/box, fill=ragLLMBg, minimum width=3.9cm, minimum height=1.5cm}}
\tikzset{rag/store/.style={rag/box, fill=ragStoreBg, minimum width=3.7cm, minimum height=1.5cm}}
\tikzset{rag/context/.style={rag/box, fill=ragContextBg, minimum width=4.6cm, minimum height=1.7cm}}
\tikzset{rag/prompt/.style={rag/box, fill=AMU_lt1!90!white, minimum width=5.2cm, minimum height=2.0cm}}
\tikzset{rag/arrow/.style={->, line width=1.1pt, >=Stealth, draw=ragArrow}}
\tikzset{rag/note/.style={font=\scriptsize, text=ragGray}}

\newcommand{\RAGVsLLMDiagram}{%
  \begin{tikzpicture}[scale=0.8, transform shape, node distance=1.5cm and 2.8cm]  
    % Vanilla prompting
    \node[rag/query] (q1) {User query};
    \node[rag/llm, right=1.8cm of q1] (llm1) {LLM only};
    \draw[rag/arrow] (q1) -- (llm1);
    \node[rag/note, below=0.15cm of q1] {Vanilla prompting};
    \node[rag/note, below=0.15cm of llm1] {Answers from training data only};

    % RAG pipeline
    \node[rag/query, below=2.0cm of q1] (q2) {User query};
    \node[rag/store, right=1.8cm of q2] (store) {Document store\\ / index};
    \node[rag/context, below=0.9cm of store] (ctx) {Retrieved chunks};
    \node[rag/llm, right=1.8cm of ctx] (llm2) {LLM with RAG};

    \draw[rag/arrow] (q2) -- (store);
    \draw[rag/arrow] (store) -- (ctx);
    \draw[rag/arrow] (q2) |- (ctx);
    \draw[rag/arrow] (ctx) -- (llm2);

    \node[rag/note, below=0.15cm of q2] {Retrieval-augmented generation};
    \node[rag/note, below=0.15cm of llm2] {Answers grounded in your data};
  \end{tikzpicture}%
}

\newcommand{\RAGAugmentedPromptDiagram}{%
  \begin{tikzpicture}[scale=0.8, transform shape, node distance=1.5cm and 2.8cm]
    \node[rag/query] (user) {User question};
    \node[rag/store, right=0.9cm of user] (index) {Hybrid retriever\\(BM25 + embeddings)};
    \node[rag/context, right=0.9cm of index] (ctx) {Top-$k$ relevant chunks};
    \node[rag/prompt, below=0.9cm of ctx] (prompt) {\textbf{Augmented prompt}\\[0.2em]
      \scriptsize System instructions + question + retrieved context};
    \node[rag/llm, left=0.9cm of prompt] (llm) {LLM answer};

    \draw[rag/arrow] (user) -- (index);
    \draw[rag/arrow] (index) -- (ctx);
    \draw[rag/arrow] (ctx) -- (prompt);
    \draw[rag/arrow] (prompt) -- (llm);
  \end{tikzpicture}%
}

\newcommand{\RAGHybridRetrievalDiagram}{%
  \begin{tikzpicture}[scale=0.8, node distance=1.4cm and 2.4cm]
    \node[rag/query] (q) {User query};

    \node[rag/store, below left=0.5cm and 1.3cm of q] (lex) {Keyword /\\BM25};
    \node[rag/store, below right=0.5cm and 1.3cm of q] (vec) {Embedding\\search};

    \node[rag/context, below=1.7cm of q, minimum width=5.0cm] (combine) {Hybrid candidate set\\(rerank, intersect, union)};
    \node[rag/store, below=0.95cm of combine, minimum width=5.0cm] (meta) {Metadata filters\\(final refinement in our lab)};

    \draw[rag/arrow] (q) -- (lex);
    \draw[rag/arrow] (q) -- (vec);
    \draw[rag/arrow] (lex) -- (combine);
    \draw[rag/arrow] (vec) -- (combine);
    \draw[rag/arrow] (combine) -- (meta);
  \end{tikzpicture}%
}

% Listing style for code snippets
\lstset{basicstyle=\ttfamily\small,keywordstyle=\color{AMU_dk2}\bfseries,commentstyle=\itshape\color{gray},showstringspaces=false,columns=flexible}

\title{Large Language Models in Data Science}
\subtitle{Week 5: Retrieval-Augmented Generation (RAG)}
\author{Sebastian Mueller}
\institute{Aix-Marseille Universit\'e}
\date{2025-2026}

\begin{document}

\begin{frame}[plain]
  \titlepage
\end{frame}

\begin{frame}{Session Overview}
  \begin{columns}[T,onlytextwidth]
    \begin{column}{0.55\linewidth}
      \textbf{Lecture (1h)}
      \begin{enumerate}
        \item From LLMs to RAG
        \item Retrieval and search basics
        \item Augmenting prompts with context
        \item Lexical, metadata, and embedding search
        \item Putting it together: RAG architectures
        \item Example: Marseille \& data science assistant
      \end{enumerate}
    \end{column}
    \begin{column}{0.4\linewidth}
      \textbf{Lab (2h)}
      \begin{itemize}
        \item Ingest a small corpus about Marseille \& AMU data science
        \item Build an embedding index with metadata
        \item Implement a simple RAG chain
        \item Compare answers with and without retrieval
        \item Optional: add citations \& reranking
      \end{itemize}
    \end{column}
  \end{columns}
\end{frame}

\section{Motivation \& Recap}

\begin{frame}{Where We Are in the Course}
  \begin{itemize}
    \item Week 1: Tokens, embeddings, transformer architecture.
    \item Week 2: Using pretrained models via Hugging Face.
    \item Week 3: Effective LLM use for coding, research, ideation.
    \item Week 4: Text classification and intent routing for an AMU chatbot.
    \item \textbf{Week 5: RAG} — combine all ingredients to build assistants grounded in your own data.
  \end{itemize}
\end{frame}

\begin{frame}{Why Do We Need RAG?}
  \begin{itemize}
    \item LLMs are trained on huge but \emph{fixed} corpora; they do not know your internal documents or freshest data.
    \item Direct prompting relies on the model ``remembering'' facts from pretraining, which leads to hallucinations and outdated answers.
    \item Many data science tasks need answers grounded in:
      \begin{itemize}
        \item institutional documents (course catalog, policies),
        \item project reports, notebooks, dashboards,
        \item domain-specific knowledge bases.
      \end{itemize}
    \item RAG attaches a \emph{retrieval system} to the LLM so answers can depend on external, updatable data.
  \end{itemize}
\end{frame}

\begin{frame}{From Vanilla LLM to RAG}
  \centering
  \RAGVsLLMDiagram
\end{frame}

\begin{frame}{From Vanilla LLM to RAG}
  \begin{itemize}
    \item Mathematically we move from $P(\text{answer} \mid \text{question})$ to
      \[
        P(\text{answer} \mid \text{question}, \text{retrieved context}).
      \]
    \item The only thing the LLM sees is the \emph{prompt}; RAG works by \textbf{augmenting the prompt} with relevant context.
  \end{itemize}
\end{frame}

\section{RAG Building Blocks}

\begin{frame}{Core Components of a RAG System}
  \begin{itemize}
    \item \textbf{Document store}: where raw texts live (internet, files, database, data lake, etc.).
    \item \textbf{Chunking}: split documents into passages that are short enough to fit in the context window.
    \item \textbf{Metadata}: structured fields such as title, source, date, language, tags.
    \item \textbf{Index}: data structures to support fast search (inverted index, vector index, or both).
    \item \textbf{Retriever}: given a query, returns relevant chunks (possibly with scores).
    \item \textbf{LLM + prompt template}: turns query + retrieved chunks into an answer.
  \end{itemize}
\end{frame}

\begin{frame}{Documents, Chunks, and Metadata}
  \begin{itemize}
    \item \textbf{Document}: a logically complete piece of text (course syllabus, FAQ page, PDF report).
    \item \textbf{Chunk}: a slice of a document (e.g., a few paragraphs) used as the atomic retrieval unit.
    \item \textbf{Metadata examples for AMU / Marseille:}
      \begin{itemize}
        \item \texttt{source\_type} (``course\_catalog'', ``housing\_info'', ``lab\_sheet'').
        \item \texttt{campus} (``Luminy'', ``Saint-Charles'', ``Aix''), \texttt{city} (``Marseille'', ``Aix-en-Provence'').
        \item \texttt{program} (``MSc Data Science''), \texttt{year}, \texttt{language}.
      \end{itemize}
  \end{itemize}
\end{frame}

\begin{frame}{Retrieval Basics: Three Signals}
  \begin{itemize}
    \item \textbf{Lexical / keyword search}
      \begin{itemize}
        \item Uses exact or fuzzy matches of words (e.g., BM25).
        \item Very precise when user vocabulary matches the documents.
      \end{itemize}
    \item \textbf{Embedding (semantic) search}
      \begin{itemize}
        \item Encode query and chunks into vectors; retrieve nearest neighbors.
        \item Captures semantics beyond exact words (``student card'' $\approx$ ``carte \'etudiant'').
      \end{itemize}
    \item \textbf{Metadata filtering}
      \begin{itemize}
        \item Restrict the candidate set using structured filters (e.g., campus = ``Marseille'').
        \item Cheap and interpretable; in our lab pipeline we apply it as a \emph{final refinement} after lexical and embedding search.
        \item In large production systems you might also use metadata as a coarse pre-filter before more expensive retrieval.
      \end{itemize}
  \end{itemize}
\end{frame}

\begin{frame}{From Bag-of-Words to BM25}
  \begin{itemize}
    \item \textbf{Bag-of-words (BoW)}:
      \begin{itemize}
        \item represent each document as a vector of word counts,
        \item simple and fast, but treats all terms and all documents equally.
      \end{itemize}
    \item \textbf{Term frequency (TF)}: words that repeat many times in a document should matter more.
    \item \textbf{Inverse document frequency (IDF)}: rare words across the corpus carry more signal than very common ones.
    \item Naive TF or TF--IDF:
      \begin{itemize}
        \item can over-favor very long documents (more opportunities to match),
        \item can over-reward repeated words without saturation.
      \end{itemize}
    \item BM25 refines this idea to give a stronger, length-aware lexical baseline.
  \end{itemize}
\end{frame}

\begin{frame}{BM25: Tunable Keyword Scoring}
  \small
  \begin{itemize}
    \item \textbf{Definition}: BM25 scores a document $d$ for a query term $t$ roughly as:
      \[
        \text{score}(t,d) \propto \text{IDF}(t)\;
        \frac{(k_1 + 1)\,\text{TF}(t,d)}{\text{TF}(t,d) + k_1\left(1 - b + b\cdot \frac{\text{document length}}{\text{average document length}}\right)}.
      \] 
      {\scriptsize Best Matching 25 was named as the 25th variant in a series of scoring
functions proposed by its creators.}
    \item \textbf{Key intuitions}:
      \begin{itemize}
        \item \emph{Term frequency saturation}: repeating a word helps, but with diminishing returns.
        \item \emph{Document length normalization}: longer documents are penalized.
        \item \emph{IDF}: rare terms across the corpus contribute more to the score.
      \end{itemize}
    \item \textbf{Tunable parameters}:
      \begin{itemize}
        \item $k_1$ controls how fast the TF effect saturates,
        \item $b$ controls how strongly document length is normalized.
      \end{itemize}
    \item In practice: a well-tuned BM25 is often the \emph{best starting point} for keyword-based retrieval in RAG systems.
  \end{itemize}
\end{frame}

\begin{frame}{Combining Keywords, and Embeddings}
  \centering
  \RAGHybridRetrievalDiagram
\end{frame}


\begin{frame}{Combining Keywords, and Embeddings}
  \begin{itemize}
    \item Typical RAG uses a \textbf{hybrid retriever}:
      \begin{itemize}
        \item run keyword search and embedding search,
        \item combine candidates by intersection / union and rerank,
        \item then apply metadata filters as a final refinement step (as in our lab).
      \end{itemize}
    \item Choice of combination controls the trade-off between recall and precision.
  \end{itemize}
\end{frame}

\section{Augmenting the Prompt}

\begin{frame}{From Retrieval to Augmented Prompt}
  \centering
  \RAGAugmentedPromptDiagram
  \vspace{0.5em}
  \begin{itemize}
    \item RAG is not a new model; it is a \emph{prompting strategy} powered by a retriever.
    \item Quality hinges on:
      \begin{itemize}
        \item what you retrieve (indexing, chunking, retriever),
        \item how you format it in the prompt (templates, instructions, citations).
      \end{itemize}
  \end{itemize}
\end{frame}

\begin{frame}[fragile]{Prompt Template for Q\&A}
  \begin{itemize}
    \item Typical template (simplified):
  \end{itemize}
  \vspace{-0.4em}
  \begin{lstlisting}[language={},basicstyle=\ttfamily\small]
You are a helpful assistant for data science students in Marseille.
Use only the information in the CONTEXT to answer.
If the answer is not in the CONTEXT, say you don't know.

CONTEXT:
{{ retrieved_chunks }}

QUESTION:
{{ user_question }}
  \end{lstlisting}
  \vspace{-0.4em}
  \begin{itemize}
    \item The retrieved chunks become part of the prompt; they \emph{augment} the original question.
    \item Clear instructions reduce hallucinations and encourage citing the provided context.
  \end{itemize}
\end{frame}

\begin{frame}{Context Size and Ranking}
  \begin{itemize}
    \item Context window is limited: we cannot paste the entire knowledge base into every prompt.
    \item We must choose:
      \begin{itemize}
        \item \textbf{Top-$k$}: how many chunks to include?
        \item \textbf{Max tokens}: truncate text to stay within the limit.
        \item \textbf{Ranking strategy}: score by similarity, recency, or custom signals (e.g., student vs. admin).
      \end{itemize}
    \item Trade-off:
      \begin{itemize}
        \item too few chunks $\Rightarrow$ model misses key information;
        \item too many chunks $\Rightarrow$ context becomes noisy and harder to use.
      \end{itemize}
  \end{itemize}
\end{frame}

\section{Example: Marseille Assistant}

\begin{frame}{Use Case: AMU / Marseille Data Science Assistant}
  \begin{itemize}
    \item Build an assistant for:
      \begin{itemize}
        \item the MSc Data Science programme in Marseille,
        \item practical life in the city for students.
      \end{itemize}
    \item Example questions:
      \begin{itemize}
        \item ``Where are the data science lectures in Marseille usually held?''
        \item ``How can I access the computer labs on Sundays?''
        \item ``What are good quiet places to study data science in Marseille?''
      \end{itemize}
    \item We want answers grounded in:
      \begin{itemize}
        \item official programme pages,
        \item campus maps and schedules,
        \item curated local tips in our own documents.
      \end{itemize}
  \end{itemize}
\end{frame}

\begin{frame}{Step 1: Ingest \& Chunk}
  \begin{itemize}
    \item Collect source documents:
      \begin{itemize}
        \item programme descriptions (PDFs, HTML),
        \item campus information and opening hours,
        \item local guides curated by teaching staff.
      \end{itemize}
    \item Normalize text (encoding, language detection, basic cleaning).
    \item Chunk into passages:
      \begin{itemize}
        \item target e.g. 200--500 tokens per chunk,
        \item keep semantic coherence (sections, headings).
      \end{itemize}
    \item Attach metadata:
      \begin{itemize}
        \item \texttt{campus=Marseille}, \texttt{program=Data Science}, \texttt{source\_url}, \texttt{section\_title}, \dots
      \end{itemize}
  \end{itemize}
\end{frame}

\begin{frame}{Step 2: Index \& Retrieve}
  \begin{itemize}
    \item Compute embeddings for each chunk using a multilingual sentence model.
    \item Store vectors in a vector index (e.g., FAISS, ScaNN, or a hosted service).
    \item At query time:
      \begin{enumerate}
        \item Parse the user question (language, intent).
        \item Run hybrid search (embeddings $+$ keywords) over all candidate chunks.
        \item Apply metadata filters as a final refinement (e.g., \texttt{campus=Marseille}, \texttt{program=Data Science}).
        \item Return top-$k$ chunks with scores after filtering.
      \end{enumerate}
    \item This ``search'' stage is independent from the LLM and can be tuned separately (retriever, metadata filters, and ranking strategy).
  \end{itemize}
\end{frame}

\begin{frame}[fragile]{Step 3: Build the Augmented Prompt - Template in code (pseudo-Python}
  \vspace{-0.6em}
  \small
  \begin{lstlisting}[language=Python]
chunks = retrieve_hybrid(query)  # BM25 + embeddings
chunks = [ # Metadata-based refinement applied after retrieval
    c for c in chunks
    if c["campus"] == "Marseille"
    and c["program"] == "Data Science"
]
context = "\n\n".join(chunk["text"] for chunk in chunks)
prompt = f"""
You are an assistant for data science students in Marseille. Answer the ques-
tion using only the CONTEXT below. If the answer is not in the CONTEXT, say you 
don't know.
CONTEXT:
{context}
QUESTION:
{query}"""
  \end{lstlisting}
\end{frame}

\begin{frame}{Step 4: Answer, Evaluate, Iterate}
  \begin{itemize}
    \item Send the augmented prompt to the chosen LLM (OpenAI, local model, etc.).
    \item Log:
      \begin{itemize}
        \item user question,
        \item retrieved chunks and scores,
        \item final answer and model parameters.
      \end{itemize}
    \item Evaluate on a small test set:
      \begin{itemize}
        \item correctness (does the answer match ground truth?),
        \item grounding (can we see the supporting chunks?),
        \item robustness (different phrasings, languages, levels of detail).
      \end{itemize}
    \item Iteratively tune:
      \begin{itemize}
        \item retriever (chunking, embeddings, ranking),
        \item prompt (instructions, formatting, citations).
      \end{itemize}
  \end{itemize}
\end{frame}

\begin{frame}{Precision and Recall Refresher}
  \begin{itemize}
    \item \textbf{Precision}: among the retrieved items, how many are relevant?
      \[
        \text{Precision} = \frac{\text{true positives}}{\text{true positives} + \text{false positives}}.
      \]
    \item \textbf{Recall}: among all relevant items in the corpus, how many did we retrieve?
      \[
        \text{Recall} = \frac{\text{true positives}}{\text{true positives} + \text{false negatives}}.
      \]
    \item Example (retrieval for a RAG system):
      \begin{itemize}
        \item there are 20 truly relevant chunks in the corpus,
        \item your retriever returns 10 chunks, 8 of which are relevant.
      \end{itemize}
      Then:
      \[
        \text{Precision} = \frac{8}{10} = 0.8,\quad
        \text{Recall} = \frac{8}{20} = 0.4.
      \]
    \item For RAG, we usually prefer \emph{high recall} at the retrieval stage, then let the LLM focus on precision.
  \end{itemize}
\end{frame}

\section{Discussion \& Takeaways}

\begin{frame}{RAG vs Fine-Tuning}
  \begin{itemize}
    \item \textbf{Fine-tuning}:
      \begin{itemize}
        \item adapts model weights,
        \item good for new behaviors or styles,
        \item expensive to update, may still hallucinate facts.
      \end{itemize}
    \item \textbf{RAG}:
      \begin{itemize}
        \item keeps base model fixed,
        \item plugs in external, updatable knowledge via retrieval,
        \item easier to iterate on (update data, retriever, or prompts).
      \end{itemize}
    \item In practice:
      \begin{itemize}
        \item start with RAG for most data-centric use cases,
        \item consider fine-tuning on top of a good RAG pipeline if behavior still needs shaping.
      \end{itemize}
  \end{itemize}
\end{frame}

\begin{frame}{Common Pitfalls}
  \begin{itemize}
    \item \textbf{Hallucinated citations}: model invents sources or mixes documents.
    \item \textbf{Shallow retrieval}: top-$k$ chunks are off-topic or redundant.
    \item \textbf{Stale data}: document store not updated as the real world evolves.
    \item \textbf{Leaky prompts}: model uses prior knowledge instead of the provided context.
    \item \textbf{Evaluation gap}: no clear test set or feedback loop.
  \end{itemize}
  \vspace{0.3em}
  \begin{itemize}
    \item Good engineering practice:
      \begin{itemize}
        \item log everything,
        \item add human review for early deployments,
        \item maintain small, curated evaluation sets.
      \end{itemize}
  \end{itemize}
\end{frame}

\begin{frame}{How the Pieces Fit Together}
  \begin{itemize}
    \item Tokens, embeddings, and transformers (Weeks 1--2) explain how LLMs represent and process text.
    \item Effective prompting (Week 3) gives us control over model behavior.
    \item Classification (Week 4) is one example of using embeddings and LLMs for structured decisions.
    \item \textbf{RAG (Week 5)} combines all of these:
      \begin{itemize}
        \item embeddings for retrieval,
        \item prompts for grounding and style,
        \item evaluation for reliability.
      \end{itemize}
    \item This is the main pattern behind many modern data science assistants.
  \end{itemize}
\end{frame}

\begin{frame}{Looking Ahead}
  \begin{itemize}
    \item In the lab you will:
      \begin{itemize}
        \item build a minimal RAG pipeline over a small Marseille / AMU corpus,
        \item experiment with keyword vs embedding vs hybrid retrieval,
        \item observe how prompt design affects grounding.
      \end{itemize}
    \item In the hackathon you can reuse these ideas:
      \begin{itemize}
        \item many impactful projects are just: ``a good RAG system + a good UI''.
      \end{itemize}
  \end{itemize}
\end{frame}

\end{document}
